// Grain — simple procedural film grain (training example)
//
// A minimal grain implementation to demonstrate the core concepts:
//   - Noise generation with per-frame temporal independence
//   - Luminance-dependent amplitude (grain is strongest in midtones)
//   - Additive blending in linear light
//
// This is a simplified version of film_grain.tex. For production work,
// use that example instead — it operates in density space and models
// cross-channel correlation for physically accurate results.

f$strength = 0.1;  // grain intensity
f$size = 1.0;     // grain size (higher = coarser)

// Noise coordinates: scale by size, offset by frame for temporal variation
float scale = 500.0 / max($size, 0.1);
float seed = fi * 97.13;
float noise = simplex(u * scale + seed, v * scale + seed * 0.7);

// Luminance-dependent amplitude: peak in midtones, taper at extremes.
// This mimics how grain visibility depends on underlying exposure.
float lum = luma(@image);
float response = 4.0 * lum * (1.0 - lum);   // parabola: 0 at black/white, 1 at mid-gray

// Apply grain additively in linear light
vec3 grain = vec3(noise * $strength * response);
@OUT = @image + grain;
